\documentclass{amsart}
% For dealing with references we use the comment environment
\usepackage{verbatim}
\newenvironment{reference}{\comment}{\endcomment}
%\newenvironment{reference}{}{}
\newenvironment{slogan}{\comment}{\endcomment}
\newenvironment{history}{\comment}{\endcomment}

% For commutative diagrams we use Xy-pic
\usepackage[all]{xy}

% We use 2cell for 2-commutative diagrams.
\xyoption{2cell}
\UseAllTwocells

% We use multicol for the list of chapters between chapters
\usepackage{multicol}

% This is generally recommended for better output
\usepackage{lmodern}
\usepackage[T1]{fontenc}

% For cross-file-references

% Package for hypertext links:
\usepackage{hyperref}

% For any local file, say "hello.tex" you want to link to please

% Theorem environments.
%
\theoremstyle{plain}
\newtheorem{theorem}[subsection]{Theorem}
\newtheorem{proposition}[subsection]{Proposition}
\newtheorem{lemma}[subsection]{Lemma}

\theoremstyle{definition}
\newtheorem{definition}[subsection]{Definition}
\newtheorem{example}[subsection]{Example}
\newtheorem{exercise}[subsection]{Exercise}
\newtheorem{situation}[subsection]{Situation}

\theoremstyle{remark}
\newtheorem{remark}[subsection]{Remark}
\newtheorem{remarks}[subsection]{Remarks}

\numberwithin{equation}{subsection}

% Macros
%
\def\lim{\mathop{\mathrm{lim}}\nolimits}
\def\colim{\mathop{\mathrm{colim}}\nolimits}
\def\Spec{\mathop{\mathrm{Spec}}}
\def\Hom{\mathop{\mathrm{Hom}}\nolimits}
\def\Ext{\mathop{\mathrm{Ext}}\nolimits}
\def\SheafHom{\mathop{\mathcal{H}\!\mathit{om}}\nolimits}
\def\SheafExt{\mathop{\mathcal{E}\!\mathit{xt}}\nolimits}
\def\Sch{\mathit{Sch}}
\def\Mor{\mathop{\mathrm{Mor}}\nolimits}
\def\Ob{\mathop{\mathrm{Ob}}\nolimits}
\def\Sh{\mathop{\mathit{Sh}}\nolimits}
\def\NL{\mathop{N\!L}\nolimits}
\def\CH{\mathop{\mathrm{CH}}\nolimits}
\def\proetale{{pro\text{-}\acute{e}tale}}
\def\etale{{\acute{e}tale}}
\def\QCoh{\mathit{QCoh}}
\def\Ker{\mathop{\mathrm{Ker}}}
\def\Im{\mathop{\mathrm{Im}}}
\def\Coker{\mathop{\mathrm{Coker}}}
\def\Coim{\mathop{\mathrm{Coim}}}

% Boxtimes
%
\DeclareMathSymbol{\boxtimes}{\mathbin}{AMSa}{"02}

%
% Macros for moduli stacks/spaces
%
\def\QCohstack{\mathcal{QC}\!\mathit{oh}}
\def\Cohstack{\mathcal{C}\!\mathit{oh}}
\def\Spacesstack{\mathcal{S}\!\mathit{paces}}
\def\Quotfunctor{\mathrm{Quot}}
\def\Hilbfunctor{\mathrm{Hilb}}
\def\Curvesstack{\mathcal{C}\!\mathit{urves}}
\def\Polarizedstack{\mathcal{P}\!\mathit{olarized}}
\def\Complexesstack{\mathcal{C}\!\mathit{omplexes}}
% \Pic is the operator that assigns to X its picard group, usage \Pic(X)
% \Picardstack_{X/B} denotes the Picard stack of X over B
% \Picardfunctor_{X/B} denotes the Picard functor of X over B
\def\Pic{\mathop{\mathrm{Pic}}\nolimits}
\def\Picardstack{\mathcal{P}\!\mathit{ic}}
\def\Picardfunctor{\mathrm{Pic}}
\def\Deformationcategory{\mathcal{D}\!\mathit{ef}}

\setlength{\emergencystretch}{3em}
\begin{document}
\title[Stacks Project, Tag 09VA]{382 --- Closed immersion pushforward detects finite-level pseudo-coherence}
\maketitle
\noindent Copyright (C) 2005--2025 Johan de Jong. Stacks Project authors. Source: \href{https://stacks.math.columbia.edu/tag/09VA}{Tag 09VA}.

\noindent Editorial difficulty note (not author text). Difficulty H1: The equivalence requires constructive local finite presentations in both directions under only pseudo-coherence of the ambient structure module. The proof uses adjunction to lift top-degree generators, compares their cones after exact pushforward and performs two separate inductions controlling the pseudo-coherence level..

\noindent Editorial provenance note (not author text). History: excerpt from revision \texttt{a0444\allowbreak 6e57e\allowbreak c1fbc\allowbreak 252a8\allowbreak 71afc\allowbreak ec775\allowbreak 2fb28\allowbreak 07b14}, perfect.tex, lines 2925--3004. Reference presentation was adapted; original-author context and core support blocks were retained or added. The mathematical wording is unchanged. Licensed under GNU FDL 1.2 or later; no invariant sections or cover texts. See ../licenses/COPYING. Context blocks reproduce author definitions and supporting results; they are not additional counted problems.

\begin{lemma}
\label{lemma-closed-push-pseudo-coherent}
Let $i : Z \to X$ be a morphism of ringed spaces such that
$i$ is a closed immersion of underlying topological spaces and such that
$i_*\mathcal{O}_Z$ is pseudo-coherent as an $\mathcal{O}_X$-module.
Let $E \in D(\mathcal{O}_Z)$. Then $E$ is $m$-pseudo-coherent
if and only if $Ri_*E$ is $m$-pseudo-coherent.
\end{lemma}

\begin{proof}
Throughout this proof we will use that $i_*$ is an exact functor, and
hence that $Ri_* = i_*$, see Modules, Lemma \href{https://stacks.math.columbia.edu/tag/01AX}{lemma i star exact (Tag 01AX)}.

\medskip\noindent
Assume $E$ is $m$-pseudo-coherent. Let $x \in X$. We will find a neighbourhood
of $x$ such that $i_*E$ is $m$-pseudo-coherent on it. If $x \not \in Z$
then this is clear. Thus we may assume $x \in Z$. We will use
that $U \cap Z$ for $x \in U \subset X$ open form a fundamental system of
neighbourhoods of $x$ in $Z$. After shrinking $X$ we may assume $E$ is
bounded above. We will argue by induction on
the largest integer $p$ such that $H^p(E)$ is nonzero. If $p < m$, then
there is nothing to prove. If $p \geq m$, then $H^p(E)$ is an
$\mathcal{O}_Z$-module of finite type, see
Cohomology, Lemma \href{https://stacks.math.columbia.edu/tag/08DN}{lemma finite cohomology (Tag 08DN)}.
Thus we may choose, after shrinking $X$, a map
$\mathcal{O}_Z^{\oplus n}[-p] \to E$ which induces a surjection
$\mathcal{O}_Z^{\oplus n} \to H^p(E)$. Choose a distinguished triangle
$$
\mathcal{O}_Z^{\oplus n}[-p] \to E \to C \to \mathcal{O}_Z^{\oplus n}[-p + 1]
$$
We see that $H^j(C) = 0$ for $j \geq p$ and that $C$ is $m$-pseudo-coherent
by Cohomology, Lemma \href{https://stacks.math.columbia.edu/tag/08CD}{lemma cone pseudo coherent (Tag 08CD)}.
By induction we see that $i_*C$ is $m$-pseudo-coherent on $X$.
Since $i_*\mathcal{O}_Z$ is $m$-pseudo-coherent on $X$ as well, we conclude
from the distinguished triangle
$$
i_*\mathcal{O}_Z^{\oplus n}[-p] \to i_*E \to i_*C \to
i_*\mathcal{O}_Z^{\oplus n}[-p + 1]
$$
and 
Cohomology, Lemma \href{https://stacks.math.columbia.edu/tag/08CD}{lemma cone pseudo coherent (Tag 08CD)}
that $i_*E$ is $m$-pseudo-coherent.

\medskip\noindent
Assume that $i_*E$ is $m$-pseudo-coherent. Let $z \in Z$.
We will find a neighbourhood of $z$ such that $E$
is $m$-pseudo-coherent on it. We will use
that $U \cap Z$ for $z \in U \subset X$ open form a fundamental system of
neighbourhoods of $z$ in $Z$. After shrinking $X$ we may assume $i_*E$
and hence $E$ is bounded above. We will argue by induction on
the largest integer $p$ such that $H^p(E)$ is nonzero. If $p < m$, then
there is nothing to prove. If $p \geq m$, then $H^p(i_*E) = i_*H^p(E)$
is an $\mathcal{O}_X$-module of finite type, see
Cohomology, Lemma \href{https://stacks.math.columbia.edu/tag/08DN}{lemma finite cohomology (Tag 08DN)}.
Choose a complex $\mathcal{E}^\bullet$ of $\mathcal{O}_Z$-modules
representing $E$. We may choose, after shrinking $X$,
a map $\alpha : \mathcal{O}_X^{\oplus n}[-p] \to i_*\mathcal{E}^\bullet$
which induces a surjection
$\mathcal{O}_X^{\oplus n} \to i_*H^p(\mathcal{E}^\bullet)$.
By adjunction we find a map
$\alpha : \mathcal{O}_Z^{\oplus n}[-p] \to \mathcal{E}^\bullet$
which induces a surjection
$\mathcal{O}_Z^{\oplus n} \to H^p(\mathcal{E}^\bullet)$.
Choose a distinguished triangle
$$
\mathcal{O}_Z^{\oplus n}[-p] \to E \to C \to \mathcal{O}_Z^{\oplus n}[-p + 1]
$$
We see that $H^j(C) = 0$ for $j \geq p$. From the distinguished triangle
$$
i_*\mathcal{O}_Z^{\oplus n}[-p] \to i_*E \to i_*C \to
i_*\mathcal{O}_Z^{\oplus n}[-p + 1]
$$
the fact that $i_*\mathcal{O}_Z$ is pseudo-coherent
and 
Cohomology, Lemma \href{https://stacks.math.columbia.edu/tag/08CD}{lemma cone pseudo coherent (Tag 08CD)}
we conclude that $i_*C$ is $m$-pseudo-coherent.
By induction we conclude that $C$ is $m$-pseudo-coherent.
By Cohomology, Lemma \href{https://stacks.math.columbia.edu/tag/08CD}{lemma cone pseudo coherent (Tag 08CD)}
again we conclude that $E$ is $m$-pseudo-coherent.
\end{proof}
\end{document}
